\uuid{EdxS}
\titre{Équation}
\auteur{bodin ; hanani ; mzari}
\organisation{exo7}
\licence{CC BY-NC-SA 4.0}
\datecreate{2020-01-20}
\niveau{L1}
\module{Analyse}
\chapitre{Fonctions circulaires et hyperboliques inverses}
\sousChapitre{Autre}
\difficulte{2}
\theme{Fonctions usuelles, équation}
\exo7qcmid{lille-1.fonctions-usuelles.012}

\contenu{
\begin{question}
\qclassification{126.00}
\qtags{motcle=équation}
Soit $(E)$ l'équation : $ 4^x-3^x=3^{x+1}- 2^{2x+1}$. Quelles sont les assertions vraies ?

\begin{answers}
    \good{$(E)$ est définie sur $\Rr$.}

    \good{$(E)$ admet une unique solution $x=1$.}

    \bad{$(E)$ admet deux solutions distinctes.}

    \bad{$(E)$ n'admet pas de solution.}
\end{answers}
\begin{explanations}
On a : $(E) \Leftrightarrow 2^{2x}+ 2^{2x+1} = 3^{x+1}+3^x \Leftrightarrow (1+2)2^{2x}=(1+3)3^{x}  \Leftrightarrow 2^{2x-2} =3^{x-1} \Leftrightarrow (2x-2)\ln 2=(x-1) \ln 3 \Leftrightarrow  x=1$.
\end{explanations}
\end{question}
}
